\documentclass[11pt]{article}
\usepackage[margin=2.4cm]{geometry}
\usepackage{booktabs}
\usepackage{graphicx}
\usepackage{amsmath}
\usepackage{times}
\usepackage[hidelinks]{hyperref}
\usepackage{caption}
\usepackage{subcaption}
\captionsetup{font=small,labelfont=bf}
\usepackage{url}
\urlstyle{same}
\setlength{\parskip}{2pt}

\title{\bf Four Ways a Retrieval-Augmented Generation Number Is Not What It Seems:\\
\large An Audit of Templates, Instructions, Pools and Tie Conventions}

\author{Fangqi Liu\textsuperscript{1,*}\quad Yuhe Liu\textsuperscript{1}\\[4pt]
\small \textsuperscript{1}School of Computer Science and Engineering, Central South University,
Changsha, China\\
\small \textsuperscript{*}\textbf{Corresponding author:} Fangqi Liu,
\texttt{fangqiliu404@gmail.com}}

\date{}

\begin{document}
\maketitle

\begin{abstract}
\noindent
A retrieval-augmented generation (RAG) result is usually reported as a single number. We audit
four choices that enter that number before any method does, and measure each in a controlled
setting. \textbf{(i) Template.} Running the default question-answering prompts shipped by
LangChain, LlamaIndex and FlashRAG verbatim on the same reader and documents yields exact
matches of 46.4, 0.0 and 0.0 on HotpotQA --- a \textbf{46-point} spread --- while a
format-insensitive containment measure puts the same three within five points (McNemar
$p=0.54$). The gap is response length: 19 versus 134 characters. \textbf{(ii) Instruction.} One
sentence inside LangChain's default drives empty-context abstention to \textbf{79--95\%},
against \textbf{0--2\%} when that sentence alone is deleted; the deletion also removes 6--8
points of the apparent benefit of retrieval. A Llama-3.2-1B reader abstains on \textbf{91.5\%}
of questions \emph{with the documents present} under that template and 0.0\% under
LlamaIndex's ($p=1.6\times10^{-55}$). \textbf{(iii) Pool.} In the ten-paragraph distractor pool
that benchmarks supply, the gold paragraph is present for 97.4\% of questions, so oracle
selection is worth only \textbf{+2.0} EM; on a 4{,}930-paragraph open corpus, full gold
coverage falls to \textbf{59.0\%} for BM25 and reranking then buys \textbf{+5.5} EM. Whether
``selection helps'' is a property of the pool. \textbf{(iv) Metric.} When confidence has few
levels, the ranking score depends on the tie convention; the difference between conventions is
exactly $\tfrac12\sum_i w_i\,\Delta r_i$, equal to $(1-\mathrm{acc})/(2L)$ for equal-width
groups, and reaches 0.19 for three-sample voting. We give a checklist that makes each of the
four explicit, and anchor the closed-book condition against FlashRAG's published setting
(F1 24.9 versus 28.4 reported).

\vspace{4pt}
\noindent\textbf{Keywords:} retrieval-augmented generation; evaluation methodology; prompt
templates; abstention; retrieval coverage; rank metrics; reproducibility

\end{abstract}

\begin{figure}[t]
\centering
\includegraphics[width=0.95\linewidth]{fig_ai_main.png}
\caption{Overview of the four inputs audited here, drawn for the two largest. Left: the same
reader, documents and questions under three published default templates, with the 46-point
exact-match spread. Right: abstention with an empty context, with and without the single
sentence that LangChain's default carries. Values are the measured ones from
Tables~\ref{tab:format} and~\ref{tab:clause}.}
\label{fig:overview}
\end{figure}
\section{Introduction}

A retrieval-augmented generation system is assembled from a retriever, a reranker, a reader and
a prompt that glues them together~\cite{lewis2020rag,gao2023survey}. Papers describe the first
three in detail and report one number for the whole. We ask what that number actually contains.

We isolate four inputs that are chosen before any method is:

\begin{enumerate}
\item \textbf{the template} --- the wording the framework ships and the practitioner copies;
\item \textbf{the instruction} --- sentences inside that template that change behaviour;
\item \textbf{the pool} --- the documents the retriever is allowed to choose from;
\item \textbf{the metric} --- how ties in a confidence score are broken.
\end{enumerate}

Each is measured on the same stack (local quantised readers, greedy decoding, BM25 over the
supplied pool), so that the four effects are comparable and their magnitudes can be placed side
by side. The result is not that any one of them is surprising in isolation --- template
sensitivity~\cite{sclar2024quantifying,voronov2024mind,mizrahi2024state,hua2025flaw} and
abstention-as-instruction~\cite{abstention2025} are both documented. The result is their
\emph{relative size}. The template is worth 46 points, the instruction 79--95 points of
abstention, the pool 2 points of oracle headroom against 5.5 points of reranking gain, and the
tie convention up to 0.19 AURC. A method that improves EM by two points is, on this stack,
smaller than three of the four.

\subsection*{Contributions}
\begin{enumerate}
\item A controlled measurement of the template effect on shipped defaults, with a
format-insensitive control that shows the readers know the same answers
(Section~\ref{sec:template}).
\item A byte-level ablation of one instruction, with three length-matched controls, showing it
controls the closed-book baseline and inflates the measured retrieval benefit
(Section~\ref{sec:clause}).
\item A direct measurement of retrieval headroom: the same reader on a supplied pool versus an
open corpus, showing that ``selection does not help'' is a property of the pool
(Section~\ref{sec:pool}).
\item A closed form for the tie-convention component of a ranking score, and a demonstration
that it is large for the confidence signals practitioners actually use
(Section~\ref{sec:metric}).
\item An anchor of the closed-book condition against a published evaluation setting
(Section~\ref{sec:anchor}).
\end{enumerate}

We claim no new method. Every component is standard. The contribution is the audit, and it
stands or falls on reproducibility: every number here is regenerated from raw per-question
predictions by a verification script.

\section{Related work}

\paragraph{Prompt and format sensitivity.}
Sclar et al.~\cite{sclar2024quantifying} show that spurious prompt-formatting features move
benchmark accuracy by tens of points; Voronov et al.~\cite{voronov2024mind} repeat the finding
for in-context learning; Mizrahi et al.~\cite{mizrahi2024state} argue for multi-prompt
evaluation. Webson and Pavlick~\cite{webson2022prompt} ask whether models use prompt meaning or
surface form. We measure the same phenomenon on the defaults frameworks ship, which is what a
practitioner actually copies, and add a containment control that separates ``the model cannot
answer'' from ``the metric cannot see the answer''.

\paragraph{Evaluation as an artefact source.}
Hua et al.~\cite{hua2025flaw} trace much reported prompt sensitivity to rigid answer matching
and show it shrinks under LLM-as-a-judge~\cite{zheng2023judging} scoring. Harness-level
reproducibility has been examined at
scale~\cite{liang2023helm,biderman2024lessons}. Our first study is a concrete instance of this
thesis with a measured magnitude; our second study is not, because deleting an instruction
changes behaviour rather than scoring.

\paragraph{Abstention.}
Willingness to abstain is largely installed by the
prompt~\cite{abstention2025}, and models can report their own uncertainty when asked
appropriately~\cite{kadavath2022language,kuhn2023semantic,farquhar2024detecting}. We extend the
instructional account to the \emph{closed-book baseline}, the standard evidence that retrieval
helps.

\paragraph{Retrieval and reranking.}
RAG spans dense retrieval~\cite{karpukhin2020dense}, fusion-in-decoder
readers~\cite{izacard2021leveraging} and adaptive variants~\cite{asai2024selfrag,jiang2023flare}.
Reranking is a standard second stage. We show that its measured gain depends on how much
headroom the pool leaves, and quantify both ends.

\paragraph{Metrics.}
Exact match and token F1 come from extractive QA~\cite{rajpurkar2016squad}. Selective
prediction and risk--coverage analysis underlie the ranking scores we audit. We give an exact
decomposition of the tie component rather than an empirical estimate.

\paragraph{Benchmarks and toolkits.}
HotpotQA~\cite{yang2018hotpotqa}, 2WikiMultiHopQA~\cite{ho2020} and
MuSiQue~\cite{trivedi2022musique}; readers from Qwen~\cite{qwen2024} and Llama
3~\cite{llama3}; templates from LangChain~\cite{langchain}, LlamaIndex~\cite{llamaindex} and
FlashRAG~\cite{flashrag2025}.

\section{Setup}
\label{sec:setup}

\paragraph{Readers.} Four instruction-tuned models, quantised to 4 bits, served locally:
Llama-3.2-1B, Qwen2.5-3B (primary), Qwen2.5-7B, Llama3-8B. Greedy decoding, fixed seed, 32-token
cap.

\paragraph{Data.} HotpotQA (distractor), 2WikiMultiHopQA and MuSiQue held-out slices, 500
questions each. For Studies A and B each question carries its ten candidate paragraphs; for
Study C we additionally build a 4{,}930-paragraph corpus from the union of all questions'
paragraphs.

\paragraph{Metrics.} Exact match after normalisation; token F1; \emph{containment} (the gold
string appears anywhere in the response); \emph{abstention} (empty or refusal). Containment is
reported wherever a template elicits prose, because EM is then uninformative by construction.

\paragraph{Scale.} 21{,}000 reader calls for Studies A and B, 2{,}000 for Study C, and the
metric results are closed-form plus a re-analysis of stored confidence scores. All predictions
are kept per question; a single script recomputes every quoted value.

\section{Study A: the template}
\label{sec:template}

The three templates are reproduced verbatim in Table~\ref{tab:templates}.

\begin{table}[t]
\centering
\caption{Default QA templates, verbatim, with provenance.}
\label{tab:templates}
\small
\begin{tabular}{@{}p{1.0cm}p{6.6cm}p{1.5cm}p{2.6cm}@{}}
\toprule
id & source & clause & output cue \\
\midrule
LC & \texttt{langchain.chains} \texttt{retrieval\_qa.prompt} & \textbf{yes} & \texttt{Helpful Answer:} \\
LI & \texttt{llama\_index} \texttt{DEFAULT\_TEXT\_QA\_PROMPT} & no & \texttt{Answer:} \\
FR & FlashRAG \texttt{base\_prompt} & no & \texttt{Answer:} \\
\bottomrule
\end{tabular}
\end{table}

\begin{figure}[t]
\centering
\includegraphics[width=0.92\linewidth]{fig_flat.png}
\caption{The same reader (Qwen2.5-3B), the same ten documents and the same 500 questions, under
three published defaults. Left: exact match spans 46 points. Right: containment, which does not
depend on answer format, is flat.}
\label{fig:flat}
\end{figure}

Figure~\ref{fig:flat} is the headline: 46.4 against 0.0 and 0.0, with containment 50.6, 47.8 and
52.2. Table~\ref{tab:format} gives all three datasets.

\begin{table}[t]
\centering
\caption{Study A, retrieval condition, $n=500$ per cell.}
\label{tab:format}
\begin{tabular}{@{}llrrrr@{}}
\toprule
dataset & template & EM & F1 & containment & abstention \\
\midrule
HotpotQA & FR & \textbf{46.4} & 58.2 & 50.6 & 0.0 \\
         & LI & 0.0 & 12.3 & 47.8 & 0.0 \\
         & LC & 0.0 & 13.3 & 52.2 & 0.8 \\
\midrule
2Wiki    & FR & \textbf{37.8} & 44.4 & 44.4 & 0.2 \\
         & LI & 0.0 & 10.2 & 40.4 & 0.2 \\
         & LC & 0.0 & 11.8 & 43.0 & 6.2 \\
\midrule
MuSiQue  & FR & \textbf{12.2} & 19.9 & 15.2 & 1.0 \\
         & LI & 0.0 & 4.5 & 10.4 & 0.0 \\
         & LC & 0.0 & 4.5 & 10.2 & 6.0 \\
\bottomrule
\end{tabular}
\end{table}

Paired comparisons on containment: FlashRAG versus LangChain $-1.6$, exact McNemar $p=0.54$;
LlamaIndex versus LangChain $-4.4$, $p=0.008$. The knowledge is the same to within five points.

The mechanism is length (Table~\ref{tab:length}). The FlashRAG template yields 19-character
answers, the others 133--141.

\begin{table}[t]
\centering
\caption{Response length and score, HotpotQA, retrieval condition.}
\label{tab:length}
\begin{tabular}{@{}lrrrr@{}}
\toprule
template & mean chars & median & EM & containment \\
\midrule
FR & \textbf{19} & 12 & 46.4 & 50.6 \\
LI & 134 & 136 & 0.0 & 47.8 \\
LC & 133 & 134 & 0.0 & 52.2 \\
LC\_minus & 139 & 140 & 0.0 & 48.4 \\
LC\_neutral & 141 & 142 & 0.0 & 45.4 \\
\bottomrule
\end{tabular}
\end{table}

A natural objection is that the short answer is present at the front of the prose and merely
needs extracting. It is not (Table~\ref{tab:trunc}): truncating to the first $w$ words recovers
at most 1.4 points.

\begin{table}[t]
\centering
\caption{Exact match after truncating to the first $w$ words.}
\label{tab:trunc}
\begin{tabular}{@{}lrrrrr@{}}
\toprule
template & full & $w=3$ & $w=5$ & $w=8$ & $w=15$ \\
\midrule
FR & 46.4 & 41.8 & 45.8 & 46.6 & 46.6 \\
LI & 0.0 & 1.4 & 0.2 & 0.0 & 0.0 \\
LC & 0.0 & 1.2 & 0.2 & 0.0 & 0.0 \\
LC\_minus & 0.0 & 0.8 & 0.2 & 0.0 & 0.0 \\
\bottomrule
\end{tabular}
\end{table}

The model mentions the gold string in about half of the LangChain responses but never leads with
it: the template changes the task from extraction to commentary, and a commentary scores zero
under exact match.

\section{Study B: the instruction}
\label{sec:clause}

LangChain's default contains one sentence that the other two do not: \emph{``If you don't know
the answer, just say that you don't know, don't try to make up an answer.''} We ablate exactly
that sentence, in three ways, holding the rest of the template byte-identical.

\begin{figure}[t]
\centering
\includegraphics[width=0.78\linewidth]{fig_abstention.png}
\caption{Abstention with an empty context, Qwen2.5-3B, $n=500$ per cell. Deleting one sentence
moves abstention by 77--93 points.}
\label{fig:abst}
\end{figure}

\begin{table}[t]
\centering
\caption{Study B, empty context, Qwen2.5-3B, $n=500$.}
\label{tab:clause}
\begin{tabular}{@{}lrrrr@{}}
\toprule
dataset & LC & LC\_minus & LC\_neutral & LC\_cleared \\
\midrule
HotpotQA & \textbf{91.2} & 1.0 & 0.2 & 1.4 \\
2Wiki    & \textbf{79.0} & 0.0 & 0.2 & 0.0 \\
MuSiQue  & \textbf{95.0} & 2.0 & 0.0 & 1.6 \\
\bottomrule
\end{tabular}
\end{table}

\begin{figure}[t]
\centering
\includegraphics[width=0.74\linewidth]{fig_controls.png}
\caption{The three controls of Study B, HotpotQA, empty context. An equal-length neutral
sentence reproduces the collapse, so the effect is content, not length.}
\label{fig:controls}
\end{figure}

The equal-length neutral replacement (217 characters against the original 209) reproduces the
collapse at 0.2\%, which rules out length as the cause (Figure~\ref{fig:controls}).

Containment in the same condition barely moves: 11.2 under LC and 14.4 under LC\_minus on
HotpotQA, 28.6 and 31.8 on 2Wiki. The sentence does not change what the reader would say; it
changes whether it says anything. Consequently it inflates the apparent benefit of retrieval by
8.2 points on HotpotQA and 6.0 on 2Wiki (Table~\ref{tab:gap}).

\begin{table}[t]
\centering
\caption{Retrieval minus closed book, in containment points.}
\label{tab:gap}
\begin{tabular}{@{}lrrrr@{}}
\toprule
dataset & LC & LC\_minus & LC\_neutral & LC\_cleared \\
\midrule
HotpotQA & 41.0 & 34.0 & 34.2 & 32.8 \\
2Wiki    & 14.4 & 8.4 & 8.2 & 11.8 \\
MuSiQue  & 10.2 & 9.2 & 8.6 & 11.0 \\
\bottomrule
\end{tabular}
\end{table}

\subsection{Reader dependence}

\begin{table}[t]
\centering
\caption{Studies A and B across four readers, HotpotQA, $n=200$ per cell unless stated.}
\label{tab:readers}
\begin{tabular}{@{}llrrr@{}}
\toprule
reader & template & EM & containment & abstention \\
\midrule
Llama3-8B & FR & \textbf{54.0} & 56.5 & 0.0 \\
          & LC & 2.0 & 45.5 & \textbf{22.0} \\
          & LI & 0.0 & 42.5 & 0.0 \\
          & LC, closed & 0.0 & 6.0 & \textbf{80.0} \\
\midrule
Qwen2.5-7B & FR & 51.5 & 57.5 & 0.0 \\
           & LC & 0.0 & 52.5 & 0.0 \\
           & LC, closed & 0.0 & 17.0 & 15.0 \\
\midrule
Qwen2.5-3B & FR & 46.4 & 50.6 & 0.0 \\
           & LC & 0.0 & 52.2 & 0.8 \\
           & LC, closed & 0.0 & 11.2 & \textbf{91.2} \\
\midrule
Llama-3.2-1B & FR & 31.0 & 33.0 & 0.0 \\
             & LC & 0.0 & \textbf{6.5} & \textbf{91.5} \\
             & LI & 0.5 & 44.5 & 0.0 \\
\bottomrule
\end{tabular}
\end{table}

The template effect is universal. Compliance is not monotone in size: Qwen2.5-7B abstains on
0.0\% of questions with documents present while Llama3-8B abstains on 22.0\%; the smallest
reader, Llama-3.2-1B, refuses 91.5\% \emph{with the documents present} under LangChain's
template against 0.0\% under LlamaIndex's (exact McNemar $p=1.6\times10^{-55}$, 183 discordant
pairs, all one direction; containment 6.5 versus 44.5).

\section{Study C: the pool}
\label{sec:pool}

Benchmarks hand the reader a pool that contains the gold paragraph by construction. In
HotpotQA-distractor the gold is present for 97.4\% of questions, so a selector can at best
reorder evidence the reader has already seen. We measured the ceiling directly: on 200 questions
where reading all ten yields 46.5 EM, an oracle that picks the best two paragraphs yields 48.5
--- a headroom of \textbf{+2.0}.

That number is a property of the pool, not of selection. We rebuilt the corpus from the union of
all 500 questions' paragraphs (4{,}930 paragraphs) and asked whether the retriever can recover
the gold paragraphs at all (Figure~\ref{fig:retrieval}).

\begin{figure}[t]
\centering
\includegraphics[width=0.95\linewidth]{fig_retrieval.png}
\caption{Study C. Left: fraction of questions whose full gold paragraph set appears in the top
ten, supplied pool versus open corpus. Right: end-to-end exact match with and without
cross-encoder reranking on the open corpus ($n=200$).}
\label{fig:retrieval}
\end{figure}

Full gold-set coverage at $k=10$ is \textbf{59.0\%} for BM25, 62.0\% for dense retrieval, 73.0\%
for BM25 plus dense and \textbf{81.0\%} for BM25 with cross-encoder reranking (coverage at
$k=20$ and $k=50$ is 74\% and 82\% for the full set). End to end, BM25 top-10 gives 36.5 EM and
cross-encoder reranking gives 42.0 --- \textbf{+5.5} points, McNemar $p=0.052$ with 19 versus 8
discordant pairs.

The two studies are consistent and the moral is quantitative: where the pool guarantees the
gold (the benchmark setting), selection has a ceiling of about two points; where it does not,
reranking is worth more than five. A paper reporting ``selection did not help'' and a paper
reporting ``reranking helped'' can both be right, about different pools.

\section{Study D: the metric}
\label{sec:metric}

The last input is not a choice of content but of arithmetic. Ranking metrics such as AURC sort
items by confidence and integrate; when confidence takes few distinct values, many items tie,
and the integration rule decides where the mass goes. This is not a small effect and it is not
a matter of taste: it has a closed form.

Let the confidence values group the $n$ items into tie groups $g=1..G$ of width
$w_g = n_g/n$, with error rate $r_g$, sorted by decreasing confidence. Then the difference
between integrating with right-endpoint rectangles and with trapezoids is
\begin{equation}
\Delta \;=\; \tfrac12 \sum_{g=1}^{G} w_g \,\Delta r_g ,
\end{equation}
where $\Delta r_g = r_g - r_{g-1}$ is the jump in cumulative error at group $g$. For
equal-width groups this becomes
\begin{equation}
\Delta \;=\; \frac{1-\mathrm{acc}}{2L},
\end{equation}
with $L$ the number of confidence levels. We verified the identity numerically on three
datasets (measured ratio 1.000) and the corollary to four decimal places.

The practical consequence is that the same predictions receive different scores depending on a
convention that papers rarely state, and that the discrepancy is largest exactly where
confidence is coarsest. Three-sample majority voting has few effective levels: at Qwen2.5-3B the
tie component reaches \textbf{0.19} AURC, and the largest tie group contains 77\% of the
questions. Token log-probabilities, being nearly continuous, are far less affected.

\section{Anchor}
\label{sec:anchor}

An audit is only interpretable if its instrument is calibrated. FlashRAG publishes its reader
configuration --- Llama3-8B-instruct, e5-base-v2, five documents --- with HotpotQA F1 of 35.3 for
standard RAG and 28.4 without retrieval.

\begin{table}[t]
\centering
\caption{Anchor against a published setting.}
\label{tab:anchor}
\begin{tabular}{@{}llrr@{}}
\toprule
condition & retriever & ours (F1) & FlashRAG reported \\
\midrule
closed book & none & \textbf{24.9} & \textbf{28.4} \\
retrieval & ours: 10-doc pool & 67.6 & 35.3 (e5, $10^6$ corpus) \\
\bottomrule
\end{tabular}
\end{table}

The closed-book condition agrees to within 3.5 points, and it is the condition with no
retriever, hence nothing else that could differ. The retrieval condition does not agree, for the
reason Study C quantifies: their retriever searches a million documents, ours a pool that
contains the answer 97.4\% of the time.

\section{Discussion}

\paragraph{Relative size.} On this stack the four effects are not commensurable with the
differences methods usually claim. Template: 46 EM points. Instruction: 77--93 points of
abstention and 6--8 points of apparent retrieval gain. Pool: 2 points of oracle headroom in the
supplied setting against 5.5 points of reranking gain in the open setting. Metric: up to 0.19
AURC. A method paper reporting a two-point EM improvement is, here, reporting something smaller
than three of the four.

\paragraph{Checklist.}
\begin{enumerate}
\item Report the template verbatim, with its source.
\item Report abstention separately from accuracy; $\mathrm{EM}=(1-u)\,c$ makes the policy term
explicit.
\item Use a short-answer cue if EM is the metric, and report containment or F1 otherwise.
\item State the pool: how often it contains the gold, and at what $k$.
\item State the tie convention, or use a metric that does not depend on it.
\end{enumerate}

\section{Threats to validity}

\emph{Scale.} Two model families, four sizes, 4-bit quantisation, one machine.

\emph{Stack.} One sparse retriever and one cross-encoder; the absolute numbers are specific to
this stack. All comparisons are within-stack.

\emph{Templates.} Three templates is a sample; other frameworks may ship others. The clause
ablation exists only for LangChain because the other two have no clause to delete.

\emph{Generality.} English multi-hop QA. Single-hop, other languages and API-based readers are
untested.

\section{Conclusion}

Four inputs enter a RAG number before any method does, and each is larger than the differences
methods usually report: a template worth 46 points, an instruction worth 77--93 points of
abstention, a pool worth 2 points in the benchmark setting and 5.5 in the open one, and a tie
convention worth up to 0.19 AURC. None requires a new method to fix; all require reporting what
was actually done.

\section*{Declarations}

\noindent\textbf{Funding.} Not applicable.

\noindent\textbf{Use of AI tools.} Figure~\ref{fig:overview} was rendered with an AI image tool
from a specification written by the authors; the numbers it displays were transcribed from the
measured results and checked against Tables~\ref{tab:format} and~\ref{tab:clause}. No other part
of the manuscript was produced with generative AI. All reported values come from the
verification script described in Section~\ref{sec:setup}.

\noindent\textbf{Competing interests.} The authors declare no competing interests.

\noindent\textbf{Ethics approval.} Not applicable; public benchmarks and locally served
open-weight models.

\noindent\textbf{Consent to participate / for publication.} Not applicable.

\noindent\textbf{Data availability.} All datasets are public
(HotpotQA~\cite{yang2018hotpotqa}, 2WikiMultiHopQA~\cite{ho2020},
MuSiQue~\cite{trivedi2022musique}). Per-question predictions and the scripts that regenerate
every table are in the project repository.

\noindent\textbf{Code availability.}
\url{https://github.com/sanke926/moxing-rag} --- runner, analysis, figure scripts and the
verification script that recomputes all quoted values.

\noindent\textbf{Author contributions.} Fangqi Liu: conceptualisation, methodology, software,
investigation, formal analysis, writing --- original draft. Yuhe Liu: data curation,
validation, writing --- review and editing. Both authors read and approved the final
manuscript.

\begin{thebibliography}{99}

\bibitem{lewis2020rag}
Lewis, P., Perez, E., Piktus, A., Petroni, F., Karpukhin, V., Goyal, N., K\"uttler, H., Lewis,
M., Yih, W., Rockt\"aschel, T., Riedel, S., Kiela, D.: Retrieval-augmented generation for
knowledge-intensive NLP tasks. In: Advances in Neural Information Processing Systems (NeurIPS)
(2020)

\bibitem{karpukhin2020dense}
Karpukhin, V., O\u{g}uz, B., Min, S., Lewis, P., Wu, L., Edunov, S., Chen, D., Yih, W.: Dense
passage retrieval for open-domain question answering. In: Proceedings of EMNLP (2020)

\bibitem{izacard2021leveraging}
Izacard, G., Grave, E.: Leveraging passage retrieval with generative models for open domain
question answering. In: Proceedings of EACL (2021)

\bibitem{gao2023survey}
Gao, Y., Xiong, Y., Gao, X., Jia, K., Pan, J., Bi, Y., Dai, Y., Sun, J., Wang, H.: Retrieval-
augmented generation for large language models: a survey.
\url{https://arxiv.org/abs/2312.10997} (2023)

\bibitem{asai2024selfrag}
Asai, A., Wu, Z., Wang, Y., Sil, A., Hajishirzi, H.: Self-RAG: learning to retrieve, generate,
and critique through self-reflection. In: International Conference on Learning Representations
(ICLR) (2024)

\bibitem{jiang2023flare}
Jiang, Z., Xu, F., Gao, L., Sun, Z., Liu, Q., Dwivedi-Yu, J., Yang, Y., Callan, J., Neubig, G.:
Active retrieval augmented generation. In: Proceedings of EMNLP (2023)

\bibitem{sclar2024quantifying}
Sclar, M., Choi, Y., Tsvetkov, Y., Suhr, A.: Quantifying language models' sensitivity to
spurious features in prompt design, or: how I learned to start worrying about prompt
formatting. In: International Conference on Learning Representations (ICLR) (2024)

\bibitem{voronov2024mind}
Voronov, A., Wolf, L., Ryabinin, M.: Mind your format: towards consistent evaluation of
in-context learning improvements. \url{https://arxiv.org/abs/2401.06766} (2024)

\bibitem{mizrahi2024state}
Mizrahi, M., Kaplan, G., Malkin, D., Dror, R., Shahaf, G., Stanovsky, G.: State of what art? A
call for multi-prompt LLM evaluation. Transactions of the Association for Computational
Linguistics (TACL) \textbf{12} (2024)

\bibitem{webson2022prompt}
Webson, A., Pavlick, E.: Do prompt-based models really understand the meaning of their prompts?
In: Proceedings of NAACL (2022)

\bibitem{hua2025flaw}
Hua, A., Tang, K., Gu, C., Gu, J., Wong, E., Qin, Y.: Flaw or artifact? Rethinking prompt
sensitivity in evaluating LLMs. In: Proceedings of EMNLP (2025).
\url{https://arxiv.org/abs/2509.01790}

\bibitem{zheng2023judging}
Zheng, L., Chiang, W.-L., Sheng, Y., Zhuang, S., Wu, Z., Zhuang, Y., Lin, Z., Li, Z., Li, D.,
Xing, E.P., Zhang, H., Gonzalez, J.E., Stoica, I.: Judging LLM-as-a-judge with MT-Bench and
Chatbot Arena. In: NeurIPS Datasets and Benchmarks Track (2023)

\bibitem{liang2023helm}
Liang, P., Bommasani, R., Lee, T., et al.: Holistic evaluation of language models. Transactions
on Machine Learning Research (TMLR) (2023)

\bibitem{biderman2024lessons}
Biderman, S., Skowron, A., Sutawika, L., et al.: Lessons from the trenches on reproducible
evaluation suites. \url{https://arxiv.org/abs/2405.14782} (2024)

\bibitem{abstention2025}
LLM abstention can be a prompt artifact, in addition to genuine uncertainty.
\url{https://arxiv.org/abs/2507.16199} (2025)

\bibitem{kadavath2022language}
Kadavath, S., Conerly, T., Askell, A., et al.: Language models (mostly) know what they know.
\url{https://arxiv.org/abs/2207.05221} (2022)

\bibitem{kuhn2023semantic}
Kuhn, L., Gal, Y., Farquhar, S.: Semantic uncertainty: linguistic invariances for uncertainty
estimation in natural language generation. In: ICLR (2023)

\bibitem{farquhar2024detecting}
Farquhar, S., Kossen, J., Kuhn, L., Gal, Y.: Detecting hallucinations in large language models
using semantic entropy. Nature \textbf{630}, 625--630 (2024)

\bibitem{flashrag2025}
Jin, J., Zhu, Y., Dou, Z., Dong, G., Yang, X., Zhang, C., Zhao, T., Yang, Z., Wen, J.-R.:
FlashRAG: a modular toolkit for efficient retrieval-augmented generation research. In:
Companion Proceedings of the ACM Web Conference (WWW), pp. 737--740 (2025).
\url{https://arxiv.org/abs/2405.13576}

\bibitem{langchain}
LangChain: \texttt{langchain.chains.retrieval\_qa.prompt}.
\url{https://github.com/langchain-ai/langchain}

\bibitem{llamaindex}
LlamaIndex: \texttt{DEFAULT\_TEXT\_QA\_PROMPT\_TMPL} in \texttt{llama\_index.core.prompts}.
\url{https://github.com/run-llama/llama_index}

\bibitem{yang2018hotpotqa}
Yang, Z., Qi, P., Zhang, S., Bengio, Y., Cohen, W.W., Salakhutdinov, R., Manning, C.D.:
HotpotQA: a dataset for diverse, explainable multi-hop question answering. In: Proceedings of
EMNLP, pp. 2369--2380 (2018)

\bibitem{ho2020}
Ho, X., Duong Nguyen, A.-K., Sugawara, S., Aizawa, A.: Constructing a multi-hop QA dataset for
comprehensive evaluation of reasoning steps. In: Proceedings of COLING, pp. 6609--6625 (2020)

\bibitem{trivedi2022musique}
Trivedi, H., Balasubramanian, N., Khot, T., Sabharwal, A.: MuSiQue: multihop questions via
single-hop question composition. TACL \textbf{10}, 539--554 (2022)

\bibitem{rajpurkar2016squad}
Rajpurkar, P., Zhang, J., Lopyrev, K., Liang, P.: SQuAD: 100,000+ questions for machine
comprehension of text. In: Proceedings of EMNLP (2016)

\bibitem{press2023measuring}
Press, O., Zhang, M., Min, S., Schmidt, L., Smith, N.A., Lewis, M.: Measuring and narrowing the
compositionality gap in language models. In: Findings of EMNLP (2023)

\bibitem{qwen2024}
Qwen Team: Qwen2.5 technical report. \url{https://arxiv.org/abs/2412.15115} (2024)

\bibitem{llama3}
Grattafiori, A., et al.: The Llama 3 herd of models. \url{https://arxiv.org/abs/2407.21783}
(2024)

\end{thebibliography}

\end{document}

